\section{Fixed sources for the chronological partial expansion}

We now place the sources of every partial branch in the same ordered
register spaces. The order is inherited from the original gate
occurrences, before any branch labels are chosen. This supplies the
common free inputs used in the two separated contractions in
\cite[proof of Theorem~5.2]{OpenAI2026PolynomialPEPS},
\texttt{04-compression.tex}, lines 342--434.

\begin{definition}[The retained source positions]
    \label{def:peps_partial_fixed_positions}
    \lean{TNLean.PEPS.PairEffect.SourceCircuit.partialPositions}
    \lean{TNLean.PEPS.PairEffect.SourceCircuit.partialSlots}
    \lean{TNLean.PEPS.PairEffect.SourceCircuit.partialSlotEquiv}
    \leanok
    \uses{thm:peps_partial_source_order}
    Let $\mathcal E$ be the original ordered source positions and let
    $\mathcal A$ be the affected parties. Retain the ordered sublist
    \[
        \mathcal E_{\mathcal A}
        =\{e\in\mathcal E:
        p_e\in\mathcal A\text{ or }q_e\in\mathcal A\}.
    \]
    Write $r=|\mathcal E_{\mathcal A}|$. Reading this sublist in order
    gives a bijection from $\{0,\ldots,r-1\}$ to its positions.
    The reference source list assigns to the corresponding position
    $e$ its grouped endpoints, its fixed halfspaces
    $\C^{u_e}$ and $\C^{v_e}$, and the zero vector. Thus the reference
    list specifies common register spaces. Its vectors will be
    replaced by the normalized vectors of each partial branch.
\end{definition}

\begin{theorem}[The common ordered register spaces]
    \label{thm:peps_partial_fixed_layout}
    \lean{TNLean.PEPS.PairEffect.SourceCircuit.mem_partialPositions_iff}
    \lean{TNLean.PEPS.PairEffect.SourceCircuit.partialSlotEquiv_apply}
    \lean{TNLean.PEPS.PairEffect.SourceCircuit.layout_partialSlots}
    \lean{TNLean.PEPS.PairEffect.SourceCircuit.partialSlot_spec}
    \lean{TNLean.PEPS.PairEffect.SourceCircuit.layout_partialSlots_eq}
    \leanok
    \uses{def:peps_partial_fixed_positions,thm:peps_partial_source_order}
    The source at reference index $i$ has the grouped original
    endpoints of its enumerated position $e_i$ and the halfspaces
    $\C^{u_{e_i}}$, $\C^{v_{e_i}}$. The endpoints remain distinct.
    Every partial branch has precisely this ordered source-register
    layout, independently of its labels. Membership in the reference
    list is equivalent to having an affected endpoint.
\end{theorem}

\begin{proof}
    \leanok
    \uses{thm:peps_source_circuit_incidence,thm:peps_affected_owner_equality,
        thm:peps_partial_source_order}
    Filter the fixed exhaustive list of original positions by endpoint
    incidence. Its lack of repetitions gives the stated bijection.
    At least one endpoint of every retained position is affected, so
    grouping exterior owners cannot identify its two endpoints.
    The ordered inventory formula of
    Theorem~\ref{thm:peps_partial_source_order} gives the same endpoint
    spaces for every branch.
\end{proof}

\begin{theorem}[Preparation in the fixed chronological sources]
    \label{thm:peps_partial_fixed_preparation}
    \lean{TNLean.PEPS.PairEffect.SourceCircuit.exists_partial_source_preparation}
    \leanok
    \uses{thm:peps_partial_fixed_layout,thm:peps_source_preparation}
    Let the chronological composition be allowed. For every partial
    choice $\xi$ there are unit vectors
    \[
        \eta_{\xi,i}\in
        \C^{u_{e_i}}\otimes\C^{v_{e_i}},
        \qquad 0\leq i<r,
    \]
    and an allowed composition $V_\xi$ containing no pair sources,
    such that the actual ordered source list of the partial branch
    consists exactly of these vectors at the reference positions.
    If $P_{\eta_\xi}$ prepares them and leaves the original input
    registers unchanged, then
    \[
        W_\xi=V_\xi P_{\eta_\xi}.
    \]
    The input register list of $V_\xi$ is the common source layout
    followed by the grouped original inputs. Its output register list
    is the grouped original output list. Both lists are independent
    of $\xi$. The assertion includes an empty source list and an
    empty set of partial choices.
\end{theorem}

\begin{proof}
    \leanok
    \uses{thm:peps_partial_source_inventory,thm:peps_partial_fixed_layout,
        thm:peps_source_preparation}
    Equal ordered layouts identify the two halfspaces at each source
    position. Transport the actual normalized source vectors along
    these identifications. The resulting reference list is exactly
    the branch's source inventory. The source-preparation theorem
    moves this inventory to the initial registers and gives an allowed
    remaining composition containing no sources. Its exact recovery
    identity is the displayed equality.
\end{proof}

\begin{definition}[The original label and vector at a position]
    \label{def:peps_partial_original_source_vector}
    \lean{TNLean.PEPS.PairEffect.SourceCircuit.choiceAt}
    \lean{TNLean.PEPS.PairEffect.SourceCircuit.sourceVectorAt}
    \lean{TNLean.PEPS.PairEffect.SourceCircuit.partialSlotChoice}
    \lean{TNLean.PEPS.PairEffect.SourceCircuit.partialSlotVector}
    \leanok
    \uses{def:peps_partial_source_choices,def:peps_partial_fixed_positions}
    For a touched gate occurrence $g$, a partial choice $\xi$ determines
    its original label $\xi_g\in\Xi_g$. Write $\nu_{e,\alpha}$ for
    the original vector of source position $e=(g,j)$ at local label
    $\alpha\in\Xi_g$. At the retained reference index $i$, with
    original position $e_i=(g_i,j_i)$, define
    \[
        \eta^{\mathrm{orig}}_{\xi,i}
        :=\nu_{e_i,\xi_{g_i}}.
    \]
    Both sides belong to the fixed coordinate space
    $\C^{u_{e_i}}\otimes\C^{v_{e_i}}$.
\end{definition}

\begin{theorem}[Recovery of the locally selected source vectors]
    \label{thm:peps_partial_original_source_vector}
    \lean{TNLean.PEPS.PairEffect.SourceCircuit.sourceVectorAt_norm}
    \lean{TNLean.PEPS.PairEffect.SourceCircuit.isTouched_of_source_endpoint}
    \lean{TNLean.PEPS.PairEffect.SourceCircuit.sourceAt_eq_mapOwner_sourceVectorAt}
    \lean{TNLean.PEPS.PairEffect.SourceCircuit.eq_partialSlotVector_of_inventory_eq}
    \lean{TNLean.PEPS.PairEffect.SourceCircuit.exists_partial_source_preparation_with_original_vectors}
    \leanok
    \uses{def:peps_partial_original_source_vector,thm:peps_partial_fixed_preparation}
    Every retained source belongs to a touched gate. All vectors
    $\nu_{e,\alpha}$ have norm one. The actual source record at a
    retained position is its original endpoint pair, grouped by
    affected ownership, together with
    $\eta^{\mathrm{orig}}_{\xi,i}$ in the fixed halfspaces.

    Consequently any exact representation of the ordered source
    inventory in the reference positions uses precisely these vectors.
    In Theorem~\ref{thm:peps_partial_fixed_preparation} one may take
    $\eta_{\xi,i}=\eta^{\mathrm{orig}}_{\xi,i}$. In particular, this
    vector depends only on the local label at its original owning
    gate, independently of every other gate label.
\end{theorem}

\begin{proof}
    \leanok
    \uses{thm:peps_source_circuit_incidence,thm:peps_partial_source_order,
        thm:peps_partial_fixed_preparation}
    A retained endpoint belongs to its owning gate, so that gate is
    touched. Induction on the chronological composition recovers its
    chosen local label and original vector. At a gate the vector is
    the given normalized source vector; composition and idle registers
    preserve its occurrence. Compare the two ordered inventory lists
    at reference index $i$. Equality of their source records gives
    equality of the vectors in their common halfspaces. Substitution
    in the fixed preparation gives the final assertion.
\end{proof}

\begin{theorem}[Two contractions for the actual sources]
    \label{thm:peps_actual_source_selective_partition}
    \lean{TNLean.PEPS.PairEffect.Word.exists_selective_partition_of_word}
    \leanok
    \uses{thm:peps_source_preparation,thm:peps_selective_partition}
    Let $W$ be an allowed composition and divide its parties into two
    sides. Choose a set of free positions in its actual ordered source
    list, containing every source that crosses the division. Let
    $P_{\mathrm{fix}}$ prepare the original source vectors at the
    remaining positions. There are an allowed composition $V$
    containing no pair sources, and two allowed compositions $A,B$
    containing no pair sources, such that
    \[
        \|A\|\leq1,\qquad \|B\|\leq1,
        \qquad
        J_{\mathrm{out}}VP_{\mathrm{fix}}
        =(A\otimes B)J_{\mathrm{free}}.
    \]
    Here $J_{\mathrm{free}}$ and $J_{\mathrm{out}}$ group the input
    and output registers by side. The equality holds on the entire
    free input space. The maps $V,A,B$ can be chosen before any
    vectors are supplied at the free source positions.

    More precisely, for every family $\zeta$ of vectors at those
    positions, let $P_\zeta^{\mathrm{free}}$ prepare that family and
    let $P_{\mathrm{fill}(\zeta)}$ prepare the whole source list,
    retaining the original vectors at fixed positions. Then
    \[
        J_{\mathrm{out}}VP_{\mathrm{fill}(\zeta)}
        =(A\otimes B)J_{\mathrm{free}}
        P_\zeta^{\mathrm{free}}.
    \]
    Taking the original source vectors gives the same equality with
    $VP_{\mathrm{fill}(\zeta)}=W$.
\end{theorem}

\begin{proof}
    \leanok
    \uses{thm:peps_source_preparation,thm:peps_selective_preparation_recovery,
        thm:peps_selective_partition}
    Enumerate the actual source list in its original halfspaces.
    The source-preparation theorem gives $W=VP_\eta$ with $V$
    allowed and containing no sources. Every fixed source is internal
    to one side and its original vector has norm one. Apply selective
    factorization to $V$ and these fixed preparations. Its equality
    holds before evaluating any free source vectors, and hence
    specializes both to arbitrary $\zeta$ and to the original vectors.
\end{proof}
